\section{Теоретическая часть}

\subsection{Выпуклое множество}

\begin{definition}[Выпуклое множество]
\label{dfn:convex_set}
Множество $S$ в вещественном векторном пространстве называется выпуклым, если выполняется: 
\[
	\lambda x+(1-\lambda)y\in S, \quad \forall x,y\in S, \quad \forall \lambda\in[0,1].
\]
\end{definition}
Иными словами, вместе с любыми двумя своими точками множество содержит весь отрезок, соединяющий эти точки.

Примерами выпуклых множеств на плоскости являются прямая, полуплоскость, круг, эллипс, треугольник и выпуклый многоугольник. Невыпуклыми являются кольцо, два непересекающихся круга, область с выемкой и звёздный многоугольник: в каждом случае можно подобрать две точки множества, отрезок между которыми частично выходит за его границу.

\subsection{Выпуклые комбинации и выпуклая оболочка}

\begin{definition}[Выпуклая комбинация]
Выпуклой линейной комбинацией точек $x_1,\ldots,x_k$ называется точка
\[
x=\sum_{i=1}^{k}\lambda_i x_i,
\qquad
\sum_{i=1}^{k}\lambda_i=1,
\qquad
\lambda_i\geq 0,\quad i=\overline{1,k}.
\]
\end{definition}

Для двух точек эта формула задаёт отрезок: $x=\lambda x_1+(1-\lambda)x_2$. Для трёх точек она задаёт точки треугольника с вершинами $x_1,x_2,x_3$, включая его границу.

\begin{definition}[Выпуклая оболочка]
	Выпуклая оболочка множества точек --- это наименьшее выпуклое множество, содержащее исходное множество.
\end{definition}
 Эквивалентно, она состоит из всех выпуклых комбинаций исходных точек. Поэтому в первой задаче выпуклая оболочка случайных точек представляется выпуклым многоугольником, вершины которого являются крайними точками исходного набора.
\subsection{Пересечение выпуклых множеств}

\begin{theorem}[О пересечении выпуклых множеств]
\label{thm:intersection}
Пусть $\{S_\alpha\}_{\alpha\in I}$ --- семейство выпуклых множеств.
Тогда их пересечение
\[
    S=\bigcap_{\alpha\in I} S_\alpha
\]
является выпуклым множеством.
\end{theorem}

\begin{proof}
Пусть $x,y\in S$ и $\lambda\in[0,1]$.\\
По определению пересечения
$x,y\in S_\alpha \quad \forall \alpha\in I$.

Так как каждое множество $S_\alpha$ выпукло, то
\[
    \lambda x+(1-\lambda)y\in S_\alpha
    \qquad \forall \alpha\in I.
\]

Следовательно, точка принадлежит каждому множеству $S_\alpha$, поэтому
\[
    \lambda x+(1-\lambda)y
    \in \bigcap_{\alpha\in I} S_\alpha
    = S.
\]

$\implies S$ является выпуклым множеством.
\end{proof}

Это свойство позволяет одновременно учитывать несколько выпуклых ограничений в задачах оптимизации: область допустимых решений сохраняет удобную геометрическую структуру.

\subsection{Проверка выпуклости многоугольника}

Для последовательных вершин $A$, $B$, $C$ вычисляется ориентированное двумерное векторное произведение
\[
\left(\overrightarrow{AB}\times\overrightarrow{BC}\right)_z
=
\begin{vmatrix}
x_B-x_A & x_C-x_B\\
y_B-y_A & y_C-y_B
\end{vmatrix}
=
(x_B-x_A)(y_C-y_B)
-
(x_C-x_B)(y_B-y_A).
\]
Его знак показывает направление поворота при переходе от ребра $AB$ к
ребру $BC$. У правильно заданного выпуклого многоугольника все
ненулевые повороты имеют один знак: все повороты происходят против или
по часовой стрелке. Смена знака означает наличие
внутреннего угла, направленного наружу, и, следовательно,
невыпуклость.

Нулевое произведение соответствует коллинеарным соседним рёбрам и не меняет направление обхода. В реализованной функции такие случаи пропускаются, поэтому многоугольники с промежуточными точками на одной стороне могут считаться выпуклыми.
